\section{Evaluation protocol}
\label{sec:protocol}
The question is about a \emph{new} model, so the unit of evaluation is a held-out model, not a held-out cell. The protocol was fixed before the final runs (study plan) and is as follows.

\paragraph{Splits.}
Models are split into folds (random 10-fold repeated a few times; leave-one-model-out as a check). Two shift splits test the ``new model resembles old ones'' assumption: \emph{leave-family-out} (family = organisation for MTEB, base model for the Open LLM Leaderboard, so a family never appears on both sides) and a \emph{temporal} split (train on the earliest 4/5 of dated models, test on the newest 1/5). Everything fitted is fitted on the training models of the fold: standardisation, factor model and rank, ridge/kernel penalties, \emph{and the selection ordering}. For a held-out model the first $k$ entries of the ordering it has scores for are revealed; models with fewer than $k$ scores are skipped for that $k$.

\paragraph{Fixed targets (protocol B, primary).}
In each repetition a random 25\% of the benchmarks (34\% when $N\le7$) are designated targets and are never selected; the error is measured on the held-out models' observed target cells only. Hence every $k$ is scored on the same cells and curves over $k$ are comparable. A secondary protocol A scores all remaining observed cells (the target set shrinks with $k$) and is used as a robustness check.

\paragraph{Metrics and uncertainty.}
The primary metric is $\nmse$, the mean over target cells of $((\hat y-y)/s_j)^2$ with $s_j$ the training standard deviation, so that predicting the training mean ($k=0$) scores about 1. We also report RMSE/MAE in score points (MTEB, HELM and LiveBench scores $\times100$; Open LLM Leaderboard scores are already 0--100) and, for a subset of runs, Spearman correlation between predicted and measured scores across held-out models. All confidence intervals are 95\% percentile bootstrap intervals over \emph{held-out models} (500 resamples; model-level sums are first averaged over repetitions, folds and random-selector draws); comparisons between selectors or rules are paired on the same models and resamples. The random selector is averaged over 10 independent draws per fold (5 on the Open LLM Leaderboard). Seeds are fixed (\texttt{--seed 0}). These intervals treat held-out models as exchangeable and are not corrected for the many comparisons made; for the main MTEB and HELM-MMLU runs we also resample organisations and strict model families (Appendix~\ref{app:extra}), which widens them modestly, and the variability across fold assignments is reported per repetition. Targets are one random draw per repetition (two to three on the MTEB blocks), so conclusions about unseen targets rest on few draws.

\paragraph{When is $k$ ``enough''? (fixed in advance)}
Let $E_{\mathrm{mean}}$ be the error of the column-mean predictor and $E_{\mathrm{floor}}$ the error of the same inference rule with \emph{all} other observed benchmarks revealed (the full-information predictor, run on the same folds and targets). The gap-closed fraction is $G(k)=(E_{\mathrm{mean}}-E_k)/(E_{\mathrm{mean}}-E_{\mathrm{floor}})$ and $k^\star(g)$ is the smallest $k$ whose bootstrap lower confidence bound for $G$ is at least $g\in\{0.8,0.9\}$. Absolute tolerances ($\nmse\le0.5,0.25,0.15$; RMSE $\le3,2$ points, upper bound of the CI) are reported as a secondary criterion. $G$ can exceed 1 when the floor estimate is itself noisy.

\paragraph{Bounds and reference points.}
The full-information floor (not a method: it needs all benchmarks), the oracle selector (uses targets, optimistic), the column mean, random selection under the same rule, and where available a measurement-noise reference from repeated runs inside MTEB result files (a lower bound on irreducible error, available for only some tasks).
